\section{Introduction}

The review-revise cycle is the fundamental mechanism by which research papers improve. Yet despite its centrality, we lack empirical characterization of how paper quality evolves across review rounds. How much improvement does a typical first revision yield? Which types of feedback produce the largest quality gains? How many rounds are needed before diminishing returns set in?

These questions are difficult to study in traditional peer review because: (1) most venue reviews provide only one round of feedback, (2) revision data is rarely shared publicly, and (3) the time between rounds (months to years) introduces confounding factors.

AI-simulated expert review provides a controlled environment for studying revision dynamics. By generating multiple rounds of structured reviews from calibrated reviewer personas, we can track score trajectories with high temporal resolution and controlled experimental conditions. We emphasize that our findings characterize revision dynamics \emph{within this specific AI review system}---they provide hypotheses about revision dynamics in general, but transfer to human review settings requires empirical validation (see Section~6).

\noindent We study revision dynamics across 14 papers in three research modules, analyzing:
\begin{enumerate}
    \item \textbf{Score trajectories}: How scores evolve across 2--4 review rounds
    \item \textbf{Revision strategy effectiveness}: Which types of revisions (P1 vs.\ P2 vs.\ P3) yield the greatest improvement
    \item \textbf{Round budgeting}: How many rounds are needed for submission readiness
    \item \textbf{Convergence patterns}: When do scores plateau, and what predicts early vs.\ late convergence
\end{enumerate}

\noindent \textbf{Scope and claims.} Given our dataset (14 papers, single author, one AI review system), we present this work as a \emph{descriptive pilot study} that identifies patterns and generates testable hypotheses. We report confidence intervals throughout and explicitly note where sample size limits generalizability. The primary contributions are the analytical framework (P1/P2/P3 decomposition, convergence predictors) and the empirical patterns they reveal, rather than universal claims about revision dynamics.
